\chapter*{Glossário complementar}
\addcontentsline{toc}{chapter}{Glossário complementar}
Os termos abaixo complementam o glossário geral do MOP.

\noindent\textbf{Contratação integrada:} regime de execução em que a mesma empresa elabora os projetos e executa as obras de cada sistema.\par\smallskip
\noindent\textbf{Entidade gestora:} associação, central de associações ou consórcio público responsável pela operação e pela manutenção do sistema após a entrega.\par\smallskip
\noindent\textbf{Modelo de gestão:} arranjo institucional, jurídico e financeiro que define a entidade responsável pela operação, pela manutenção e pela cobrança dos serviços de saneamento rural.\par\smallskip
\noindent\textbf{Plano Integrado de Propriedade (PIP):} plano elaborado pelo serviço de assistência técnica e extensão rural contratado no Componente 2 para cada propriedade rural aderente ao PSH-PR, com as práticas e as intervenções recomendadas.\par\smallskip
\noindent\textbf{Sistema coletivo de abastecimento de água:} conjunto de captação, tratamento, reservação, rede de distribuição e ligações domiciliares que abastece uma comunidade.\par\smallskip
\noindent\textbf{Solução individual de esgotamento sanitário:} sistema descentralizado que coleta, trata e promove a disposição final do esgoto de um domicílio no próprio lote.\par\smallskip
\noindent\textbf{Termo de Recebimento Definitivo (TRD):} documento que formaliza o recebimento final do sistema pelo IAT, após a pré-operação e o aceite da entidade gestora.\par\smallskip
\noindent\textbf{Trabalho técnico-social:} conjunto de atividades de mobilização, cadastro, diagnóstico, educação sanitária e ambiental e organização comunitária junto às famílias.\par\smallskip
\noindent\textbf{Unidade regional de apoio:} estrutura do modelo de gestão, com equipe técnica, sede e veículo, que presta suporte técnico às associações de uma região.\par\smallskip
